% !PRINTED 337
\weberpage{337}

By (12) and (7) the congruence (15) is equivalent
to the $n^2$ congruences

\[
y_{h,k}=\overset{r}{\Sigma}\,x_{r}\,c_{h,k}^{(r)}\equiv 0\pmod{p} % Source scan prints r above Sigma and no lower index.
\]
and from (6) there then follows:
\[
x_s\equiv 0\pmod{p},
\]
and this says that the substitutions (17) are linearly
independent for every modulus $p$.

\gesperrt{According to Gauss (Disq. aritm. art. 235) a form $T$
is said to be compounded or composed of the two forms $U, V$ when $T$ passes over,
by a bilinear substitution for the variable $t$ whose equations are
linearly independent for every modulus, into the product $UV$.}

Accordingly we have the theorem:

6. \gesperrt{The form that belongs to the product of two ideals $\mathfrak{a}, \mathfrak{b}$
is compounded from the forms of the ideals $\mathfrak{a}$ and $\mathfrak{b}$}\footnote{The theorems that \gesperrt{Gauss} proves at the passage mentioned by a very
lengthy calculation for binary quadratic forms are here derived in the greatest
generality by simple considerations. Cf. \gesperrt{Dirichlet-Dedekind}, Vorlesungen \"uber Zahlentheorie, \S 182. For binary quadratic
forms: \gesperrt{Dedekind}, Crelles Journal, Bd. 129; \gesperrt{H. Weber}, G\"ottinger Nachrichten 1907.}.

\begin{center}\rule{0.2\linewidth}{0.4pt}\end{center}
